% !TeX program = lualatex
% Compile twice: lualatex 81.tex
% Standalone research notebook; original section preserved.
% Research additions: 27 September 2026. No external TeX input or BibTeX file.
\documentclass[11pt,a4paper]{article}
\usepackage[margin=24mm,headheight=14pt]{geometry}
\usepackage{fontspec}
\setmainfont{Latin Modern Roman}
\setsansfont{Latin Modern Sans}
\setmonofont{Latin Modern Mono}
\usepackage{amsmath,amssymb,mathtools}
\usepackage{unicode-math}
\setmathfont{Latin Modern Math}
\usepackage{microtype}
\usepackage{enumitem,xurl,needspace}
\usepackage[dvipsnames]{xcolor}
\usepackage{hyperref}
\hypersetup{unicode=true,colorlinks=true,linkcolor=MidnightBlue,urlcolor=MidnightBlue,
 pdftitle={Research notebook - Problem 81},pdfauthor={Open Problems Math},bookmarksnumbered=true}
\usepackage{bookmark}
\usepackage{titlesec}
\titleformat{\section}{\Large\bfseries\raggedright}{\thesection}{0.7em}{}
\usepackage{fancyhdr}
\pagestyle{fancy}\fancyhf{}
\fancyhead[L]{\small\sffamily Open Problems Math | Research notebook}
\fancyhead[R]{\small\sffamily Problem 81 | 27 September 2026}
\fancyfoot[C]{\thepage}
\setlength{\parindent}{0pt}
\setlength{\parskip}{5pt plus 1pt}
\setlength{\emergencystretch}{3em}
\setcounter{tocdepth}{1}
\setcounter{secnumdepth}{1}
\setlist[itemize]{leftmargin=3em,itemsep=5pt,topsep=4pt,parsep=0pt}
\allowdisplaybreaks
\numberwithin{equation}{section}
\newcommand{\N}{\mathbb{N}}
\newcommand{\Z}{\mathbb{Z}}
\newcommand{\Q}{\mathbb{Q}}
\newcommand{\R}{\mathbb{R}}
\newcommand{\C}{\mathbb{C}}
\newcommand{\F}{\mathbb{F}}
\newcommand{\A}{\mathbb{A}}
\newcommand{\PP}{\mathbb P}
\newcommand{\E}{\mathbb{E}}
\newcommand{\eps}{\varepsilon}
\newcommand{\dd}{\,\mathrm d}
\newcommand{\sref}[2]{\hyperlink{source:#1:#2}{[S#2]}}
\newcommand{\eref}[2]{\hyperlink{extra:#1:#2}{[E#2]}}
\newif\ifshowcatalogue
\showcataloguetrue
\newcommand{\cataloguescope}[1]{\ifshowcatalogue
 \paragraph{Exact scope in the supplied catalogue.}
 \begin{quote}\small #1\end{quote}\fi}
\begin{document}
\setcounter{section}{80}
\section{\texorpdfstring{The Lindeloef Hypothesis on the Growth of the Riemann Zeta Function}{The Lindeloef Hypothesis on the Growth of the Riemann Zeta Function}}\label{81}
\subsection{Definitions and mathematical statement}
Let $\zeta$ be the meromorphically continued Riemann zeta function. The Lindelöf assertion is
\[
 \forall\varepsilon>0\ \exists C_\varepsilon,T_\varepsilon>0\quad
 |\zeta(\tfrac12+it)|\le C_\varepsilon|t|^\varepsilon
 \quad(|t|\ge T_\varepsilon).
\]
Thus growth on the critical line is smaller than every fixed positive power, with a constant allowed to depend on that power. The statement is a pointwise growth bound, not merely a mean-value estimate.
\par\smallskip\textit{Formulation and concept references:} \sref{81}{1}.
\cataloguescope{Prove that for every epsilon>0 there is a constant C\_epsilon such that |zeta(1/2+it)| <= C\_epsilon |t|\textasciicircum{}epsilon for all sufficiently large real |t|.}
\subsection{Short English statement}
Does the zeta function on the critical line grow more slowly than every fixed positive power of height?
% BEGIN RESEARCH_ADDITION ATTEMPT_81
\subsection{Research attempt}

\paragraph{Current frontier and source status.}
Florea's 2026 survey distinguishes pointwise subconvexity from the high-moment
problem, and records the Bourgain exponent $13/84$ in the pointwise
landscape \eref{81}{1}. Any one positive exponent is weaker than Lindel\"of.
The retained Unterberger source has a withdrawal at version~48; later
versions exist, so the withdrawal is not silently attributed to every
later revision. No posted full-proof claim is used here as an established
input \eref{81}{2}. The attempt focuses on the exact quantitative relation
between high moments and individual peaks.

\paragraph{Attack route.}
One route estimates the approximate-functional-equation Dirichlet polynomials
pointwise. A second proves all fixed high-moment upper bounds and then
prevents a large value from being too narrowly localized. The latter route
is pursued with an intentionally crude but explicit derivative estimate.
The mathematical reduction is complete, while the moment bounds remain
unproved. A family of smooth spikes then tests why finitely many successful
moments cannot close the argument.

\paragraph{Attempt: a peak must occupy a definite interval.}
For $T$ large and $t\in[T/2,3T]$, one has a coarse bound
\begin{equation}\label{eq81:derivative}
 |\zeta'(1/2+it)|\leq CT^2.
\end{equation}
Here is enough justification without invoking the desired hypothesis.
Euler summation in the half-plane $\sigma>0$ writes $\zeta(s)$ as a
sum up to $N$, the term $N^{1-s}/(s-1)$, and an integral involving a
bounded sawtooth function, bounded by
$O((1+|s|)N^{-\sigma}/\sigma)$. Take $N$ of order $T^2$.
On a fixed small disk about $1/2+it$, say radius $1/4$, the partial sum
is $O(T^{3/2})$, the pole term and remainder are smaller polynomial terms,
and the disk is far from $s=1$. Cauchy's derivative estimate on that disk
gives $O(T^{3/2})$, which in particular implies \eqref{eq81:derivative}.
This deliberately loses powers but is independent of Lindel\"of.

Let $t_0\in[T,2T]$ and $M=|\zeta(1/2+it_0)|$. The same coarse bound
ensures that $w=M/(2CT^2)\leq1$ after increasing $C$.
On $|t-t_0|\leq w$, the derivative bound and the reverse triangle inequality
give $|\zeta(1/2+it)|\geq M/2$. Therefore, for every fixed integer $k\geq1$,
\begin{equation}\label{eq81:peak}
 I_k(T):=\int_{T/2}^{3T}|\zeta(1/2+it)|^{2k}\,dt
 \geq\frac{M^{2k+1}}{C\,2^{2k}T^2}.
\end{equation}
If $M=0$ the desired conclusion is immediate. Otherwise the interval has
positive length and remains inside the integration range. This establishes
the peak estimate without assuming that maxima have a fixed positive
width independent of their height.

\paragraph{Attempt: the all-moments sufficient package.}
Assume that for every fixed $k\geq1$ and $\eta>0$,
\begin{equation}\label{eq81:moments}
 I_k(T)\ll_{k,\eta}T^{1+\eta}.
\end{equation}
Combining \eqref{eq81:peak} with \eqref{eq81:moments} yields
\[
 |\zeta(1/2+it_0)|\ll_{k,\eta}T^{(3+\eta)/(2k+1)}.
\]
Given the target $\varepsilon>0$, choose $k$ once and for all large enough
that $(3+1)/(2k+1)<\varepsilon$, then use \eqref{eq81:moments} with
$\eta=1$. Constants and thresholds are allowed to depend on this fixed
choice. The result is the required bound throughout $[T,2T]$.
Dyadic covering handles all large positive $t$; complex conjugation handles
negative $t$. Thus the all-fixed-moments package proves the exact pointwise
Lindel\"of assertion.

Conversely Lindel\"of gives \eqref{eq81:moments}: for a fixed $k,\eta$,
use its pointwise estimate with exponent $\eta/(2k)$ and integrate over
an interval of length $O(T)$. Hence this moment package is equivalent to
the target, in agreement with the standard relation discussed in
\eref{81}{1}. The elementary proof above is a reconstruction, not a new
proof of the unproved moment bounds. In particular $k$ is chosen before
$T$; no unproved uniform control of constants for $k\to\infty$ is used.

\paragraph{Attempt: why a finite list of moments cannot suffice by itself.}
Fix $K\geq1$ and choose a smooth nonnegative bump $\psi$ supported in
$[-1,1]$ with $\psi(0)=1$. Put $a=3/(2K+1)>0$ and choose a rapidly
increasing sequence $T_j$, with the resulting bump supports disjoint.
Define on the positive real axis
\[
 F(t)=\sum_j T_j^a\,
       \psi\left(\frac{t-T_j}{T_j^{a-2}}\right).
\]
On each bump $|F'(t)|\ll T_j^2$. Its contribution to the $2k$th moment is
\[
 T_j^{2ka}T_j^{a-2}\int\psi(u)^{2k}\,du
       \asymp_k T_j^{(2k+1)a-2}\leq C_kT_j
          \qquad(k\leq K).
\]
Taking the $T_j$ sufficiently separated gives the corresponding
$O(T)$-type dyadic moment bounds for all $k\leq K$, while
$F(T_j)=T_j^a$ violates every subpower bound with exponent below $a$.
This supplies an explicit countermodel to an inference based only on
finitely many moments and a coarse derivative bound. It is a smooth
function model, not a zeta function, and no functional equation or Euler
product is claimed for it.

\paragraph{Stress test: the arithmetic input still absent.}
A Dirichlet polynomial of length about $T^{1/2}$ becomes one of length
about $T^{k/2}$ after taking its $k$th power. A generic mean-square estimate
with a cost comparable to interval length plus polynomial length is then
too expensive for large fixed $k$. The arithmetic correlations of its
coefficients, rather than the mere existence of a polynomial representation,
are what must improve the estimate. Treating the long coefficient sums
as orthogonal beyond the available integration resolution discards the
unproved off-diagonal problem.

The argument did not obtain \eqref{eq81:moments} for any new $k$. A theorem
about a Barnes zeta function, an average over a different $L$-function
family, or a bound holding on most ordinates cannot be substituted for
that assertion. The spike example explains why exceptional large ordinates
remain relevant even after strong mean-value information. The proof also
does not use RH, and it does not claim that Lindel\"of implies RH.

\paragraph{Outcome.}
\textbf{Outcome: Conditional reduction.}

The notebook supplies a quantitative, fully justified moment-to-pointwise
reduction and a smooth countermodel to the finite-moment shortcut. Its
all-moments hypothesis is itself an equivalent unresolved route, not a
weaker theorem already available. No improvement to pointwise zeta bounds
or new high-moment estimate is obtained. A completion requires the actual
arithmetic cancellation in \eqref{eq81:moments}, or a different direct
pointwise argument. No novelty claim is made for the classical equivalence.

% Keep the local bibliography together rather than leave a source fragment.
\needspace{170mm}
% END RESEARCH_ADDITION ATTEMPT_81
\subsection{Sources}
\begin{itemize}\raggedright
\item[\textnormal{[S1]}]\hypertarget{source:81:1}{} Alexandra Florea, A survey of moment bounds for zeta(s): From Heath-Brown's work to the present, Journal of the London Mathematical Society (2026).
\url{https://londmathsoc.onlinelibrary.wiley.com/doi/10.1112/jlms.70376}.
{\small\textit{Catalogue formulation source.}}
\item[\textnormal{[S2]}]\hypertarget{source:81:2}{} Analogues of the Lindelöf Hypothesis for the Barnes multiple zeta function and related problems.
\url{https://arxiv.org/abs/2606.02407}.
{\small\textit{Catalogue research/status reference; not a proof endorsement.}}
\item[\textnormal{[S3]}]\hypertarget{source:81:3}{} An approach to the Lindelöf Hypothesis for Dirichlet L-functions.
\url{https://arxiv.org/abs/2602.05731}.
{\small\textit{Catalogue research/status reference; not a proof endorsement.}}
\item[\textnormal{[S4]}]\hypertarget{source:81:4}{} Pseudodifferential arithmetic, Riemann and Lindelöf hypotheses.
\url{https://arxiv.org/abs/2208.12937}.
{\small\textit{Catalogue research/status reference; not a proof endorsement.}}
\item[\textnormal{[S5]}]\hypertarget{source:81:5}{} Pseudodifferential arithmetic, Riemann and Lindelöf hypotheses (withdrawn v48).
\url{https://arxiv.org/abs/2208.12937v48}.
{\small\textit{Catalogue research/status reference; not a proof endorsement.}}
% BEGIN RESEARCH_ADDITION REFERENCES_81
\item[\textnormal{[E1]}]\hypertarget{extra:81:1}{} Alexandra Florea,
\emph{A survey of moment bounds for $\zeta(s)$: From Heath-Brown's work to
the present}, Journal of the London Mathematical Society (2026),
DOI 10.1112/jlms.70376, introduction and moment/Lindel\"of discussion.
\url{https://londmathsoc.onlinelibrary.wiley.com/doi/full/10.1112/jlms.70376}.
{\small\textit{Inspected 27 September 2026; the moment-to-peak proof used
here is supplied explicitly and is not presented as a new equivalence.}}
\item[\textnormal{[E2]}]\hypertarget{extra:81:2}{} Andr\'e Unterberger,
\emph{Pseudodifferential arithmetic, Riemann and Lindel\"of hypotheses},
arXiv:2208.12937v48, version-specific withdrawal notice; compare the
submission history and later revisions.
\url{https://arxiv.org/abs/2208.12937v48}.
\url{https://arxiv.org/abs/2208.12937}.
{\small\textit{History inspected 27 September 2026. The withdrawal is
attributed to v48 only; neither later reposting nor the original title is
used as evidence of an established proof.}}
% END RESEARCH_ADDITION REFERENCES_81
\end{itemize}

\end{document}
